\section{Conclusion}
Feedback loops in small software projects fail in the interval between ``raised'' and ``acted on'', and they fail invisibly because that interval is unmeasured. \ffrs{} makes the interval visible and hands its first half to an agent. The model is deliberately small: five stages, each a timestamp recorded by the system that owns it; four metrics that follow from those timestamps; and one structural rule---the Respond stage is performed by a scheduled agent, bounded by human checkpoints (merge for code; accept-and-confirm for proposals). The reference implementation adds no system of record beyond the tracker a project already has: a 4\,KB widget, one serverless function, a private sidecar for the only personal datum, and a scheduled cloud coding agent. It was built in about a working day, is removable through a three-layer toggle, and has run in production on scaledaiops.org since 18~August~2026, where a real request has already been captured, resolved by an agent-authored pull request, reviewed in a minute, merged and closed.

We claim no more than the evidence supports: that a minimal, agentic feedback loop is \emph{sufficient} for a small project under conditions we make explicit, and that its behaviour can be measured by anyone from public data. Future work is measurement over a longer window and more sites, a second tracker and runtime (for example a Jira-backed, container-hosted instance) to demonstrate the substitutability claimed in Table~\ref{tab:roles}, closure notifications that do not depend on e-mail, and a study of reviewer effort on agent-authored changes---the quantity that ultimately decides whether the human checkpoint stays a checkpoint or becomes the new bottleneck.
